% Locked-test results. Written before any locked-test number was read (2026-11-15).
% Every table is \input from reports/final/<name>.tex, which scripts/final_reports.py writes
% (all-or-nothing, fixed order) once all twelve final kernels are collected. Until then each
% table prints a "pending" line, so the draft compiles and no number is typed by hand.
% tests/test_paper_test_results.py checks the wiring and that this file contains no digits
% outside comments, labels and file names.
\section{Locked-test results}
\label{sec:test-results}

\newcommand{\finaltable}[3]{%
  \IfFileExists{../reports/final/#1.tex}{%
    \begin{table}[t]\centering\small
    \resizebox{\linewidth}{!}{\input{../reports/final/#1.tex}}
    \caption{#2}\label{#3}
    \end{table}}{%
    \par\noindent\emph{Pending: \texttt{\detokenize{reports/final/#1.tex}} is written by
    \texttt{scripts/final\_reports.py} after all locked-test kernels are collected.}\par}}
% final_slices.tex is a longtable (it runs over several pages), so it cannot sit in a float or a
% \resizebox; it is \input directly under a run-in heading.
\newcommand{\finallongtable}[3]{%
  \IfFileExists{../reports/final/#1.tex}{%
    \par\medskip\noindent\textbf{#2}\label{#3}\par
    {\scriptsize\setlength{\tabcolsep}{4pt}\input{../reports/final/#1.tex}}}{%
    \par\noindent\emph{Pending: \texttt{\detokenize{reports/final/#1.tex}} is written by
    \texttt{scripts/final\_reports.py} after all locked-test kernels are collected.}\par}}

% The figure is drawn by scripts/final_figures.py from reports/final/final_table.json only (run it after
% final_reports.py); until its PDF exists a pending line is printed.
\newcommand{\finalfigure}[3]{%
  \IfFileExists{figures/#1.pdf}{%
    \begin{figure}[t]\centering
    \includegraphics[width=\textwidth]{figures/#1.pdf}
    \caption{#2}\label{#3}
    \end{figure}}{%
    \par\noindent\emph{Pending: \texttt{\detokenize{paper/figures/#1.pdf}} is drawn by
    \texttt{scripts/final\_figures.py} from the accepted locked-test table.}\par}}

The locked test is evaluated once per arm and training seed. The prompt, token limit,
reference arm and evaluation code were fixed before the first locked-test kernel ran, and
every analysis script below was committed before any locked-test number was read
(Section~\ref{sec:repro}). The tables below are generated by committed scripts from
per-sample correctness flags; none is edited by hand. Section~\ref{sec:analysis} fixes what
each of them reports and which one decides the hypothesis.

\paragraph{Primary statistic.}
The pre-registered primary statistic is the six-cell mean of Section~\ref{sec:setup}, pooled
over seeds, with a root-graph bootstrap interval stratified by split. Its reading (support,
no support, or inconclusive) follows the rule fixed on the development set and is computed by
the script, not chosen after looking at the table.
\finaltable{final_primary}{Primary statistic on the locked test: six-cell mean and task means
per arm (seed-pooled, 95\% interval), and each difference with its pre-registered reading and
the seed-matched differences.}{tab:test-primary}
% Generated by scripts/test_reading.py (rule reading-v1, written after the test was opened, descriptive only).
\IfFileExists{../reports/final/test_reading.tex}{% generated by scripts/test_reading.py, rule reading-v1, from reports/final/final_primary.json; do not edit
\paragraph{Reading of the primary table.} Each difference below is pooled over seeds with its 95\% root-graph interval, as in Table~\ref{tab:test-primary}; the reading is the label the pre-registered rule assigned, copied from the script output. Only the D$-$B row decides the hypothesis.
\begin{itemize}
\item \textbf{D$-$B:} six-cell mean $-0.105$ [$-0.117$, $-0.093$]; reachability $-0.005$ [$-0.020$, $+0.011$], critical edge $-0.206$ [$-0.225$, $-0.186$]; seed-matched differences from $-0.115$ to $-0.096$ (3 seeds); same sign on every split: yes; recorded reading \texttt{negative}.
\item \textbf{E$-$B:} six-cell mean $+0.065$ [$+0.048$, $+0.082$]; reachability $+0.095$ [$+0.074$, $+0.115$], critical edge $+0.036$ [$+0.011$, $+0.061$]; seed-matched differences from $+0.061$ to $+0.068$ (3 seeds); same sign on every split: yes; recorded reading \texttt{positive\_clear}.
\item \textbf{C$-$B:} six-cell mean $-0.224$ [$-0.239$, $-0.208$]; reachability $-0.132$ [$-0.151$, $-0.113$], critical edge $-0.316$ [$-0.339$, $-0.292$]; seed-matched differences from $-0.238$ to $-0.212$ (3 seeds); same sign on every split: yes; recorded reading \texttt{negative}.
\item \textbf{E$-$D:} six-cell mean $+0.171$ [$+0.154$, $+0.188$]; reachability $+0.099$ [$+0.080$, $+0.120$], critical edge $+0.242$ [$+0.216$, $+0.267$]; seed-matched differences from $+0.163$ to $+0.176$ (3 seeds); same sign on every split: yes; recorded reading \texttt{positive\_clear}.
\end{itemize}
}{}

\paragraph{Per split.}
Non-edit macro accuracy per locked split, and each arm's difference to the reference arm.
\finaltable{final_table}{Non-edit macro accuracy per locked split (seed-pooled, 95\% root-graph
interval) and difference to the reference arm.}{tab:test-splits}
\finalfigure{ood_test}{Non-edit macro accuracy on the six locked splits, the values of
the per-split table. Marker: pooled over the training seeds; bar: its root-graph
bootstrap interval; small dots: single seeds. The accuracy axis starts at zero. The splits hold
different root graphs, so the figure shows levels; a difference between two splits has no
interval and no claim rests on it. Generated by
\texttt{scripts/final\_figures.py}.}{fig:test-ood}

\paragraph{Secondary analyses.}
Twin pairs, invariance pairs, density-matched slices, answer slices, the full metric set and
the frozen-base floor are secondary (Section~\ref{sec:analysis}); no claim rests on any of
them alone.
\finaltable{final_pairs}{Answer-flipping twin pairs on the locked test: both right, one right,
both wrong.}{tab:test-pairs}
\finaltable{final_invariance}{Invariance pairs on the locked test: agreement on relabelled and
restyled copies.}{tab:test-invariance}
\finaltable{final_density}{Density-matched slice on the locked test.}{tab:test-density}
\finaltable{final_metrics}{Full metric set on the locked test (strict and lenient scores,
format failures, repair detail).}{tab:test-metrics}
\finaltable{final_floor}{Frozen-base floor (no adapter) on the locked splits.}{tab:test-floor}
\finallongtable{final_slices}{Accuracy by slice on the locked test: true answer, image variant,
family, size, layout and style, per task (summed over seeds).}{tab:test-slices}
